% !TEX root = ../jasonliu2026_Thesis.tex

\section*{Nederlandstalige Samenvatting}
\addcontentsline{toc}{section}{Nederlandstalige Samenvatting}

Geautomatiseerde testgeneratie\-tools zoals EvoSuite en Randoop produceren
unit tests met niet-informatieve identifiers zoals \texttt{var0} en
\texttt{test42}, die het begrip belemmeren en de praktische waarde van de
gegenereerde uitvoer voor de ontwikkelaars die deze moeten onderhouden
verminderen. Het hernoemen van deze identifiers naar betekenisvolle namen
is een noodzakelijke nabewerkingsstap, maar bestaande niet-LLM-benaderingen
behalen hoogstens $F_1 = 0.308$ en er bestaat geen schaalbaar raamwerk
voor leesbaarheids\-evaluatie. Token-niveau nauwkeurigheid en holistische
leesbaarheid zijn niet hetzelfde: fine-tuning verbetert $F_1$ tot $0.839$
maar verlaagt tegelijkertijd de LLM-leesbaarheids\-score, terwijl een
zero-shot proprietair model de hoogste leesbaarheid van alle varianten
behaalt ondanks een vergelijkbare $F_1$, wat aantoont dat een model kan
leren om referentielabels te reproduceren zonder namen te produceren die
leesbaarder zijn voor een ontwikkelaar. Geen van beide metrieken is daarom
op zichzelf voldoende als grondwaarheid, en beide zouden samen gerapporteerd
moeten worden bij de evaluatie van identifier-hernoem\-tools; ter
ondersteuning hiervan dragen wij \textsc{CodeReader} bij, een
configureerbaar open-source leesbaarheids\-beoordelings\-tool, en tonen
wij aan dat zelfs minimale fine-tuning van 15\% van de trainingsstappen
het grootste deel van het token-niveau voordeel oplevert.
